% TOP_PROBLEM: {"id": 167, "title": "Disjoint Sumsets Construction", "category_id": 2, "category": {"id": 2, "name": "combinatorics", "display_name": "Combinatorics", "description": "Counting problems, graph theory, discrete structures.", "slug": "combinatorics", "order_index": 2, "created_at": "2026-07-31T15:26:25.671Z"}, "status": "open"}
% FIRST_POSTED: null
% ENTRY_KIND: statement_only
% Imported from: batch03.tex; UnsolvedMath ID 167
\documentclass[11pt]{article}
\usepackage[margin=25mm]{geometry}
\usepackage{fontspec}
\setmainfont{Latin Modern Roman}
\usepackage{amsmath,amssymb,mathtools}
\usepackage{unicode-math}
\setmathfont{Latin Modern Math}
\usepackage{xurl,hyperref}
\hypersetup{colorlinks=true,urlcolor=blue}
\setcounter{secnumdepth}{0}
\setlength{\parindent}{0pt}
\setlength{\parskip}{5pt}
\setlength{\emergencystretch}{3em}
\newcommand{\sref}[2]{\hyperlink{source:#1:#2}{[S#2]}}
\newcommand{\cataloguescope}[1]{\paragraph{Recorded target.}#1}
\begin{document}
% UnsolvedMath ID 167; source catalogue code GREEN-079
\setcounter{section}{166}
\section{Disjoint Sumsets Construction}\label{167}
{\small\sffamily UnsolvedMath ID 167 · Combinatorics}
\subsection{Definitions and mathematical statement}

\paragraph{Shared notation.}$\mathbb N=\{1,2,\ldots\}$, and $\mathbb Z,\mathbb Q,\mathbb R,\mathbb C$ denote the integers, rationals, reals and complex numbers. Any other conventions are specified locally.


All groups here are finite abelian and written additively. For subsets $B,C$ of a group, $B+C=\{b+c:b\in B,c\in C\}$. Equality $|B+C|=|B||C|$ means that the addition map $B\times C\to B+C$ is injective. The asymptotic assertions below mean a sequence of constructions indexed by integers $n\to\infty$ along an unbounded subsequence. A quantity $n^{2+o(1)}$ has logarithm divided by $\log n$ tending to $2$.
\paragraph{Statement recorded in the supplied source.}
Are there arbitrarily large integers $n$, finite abelian groups $H_n$, and subsets $A_1,\ldots,A_n,B_1,\ldots,B_n\subset H_n$ such that
\[
 |H_n|=n^{2+o(1)},\qquad
 \min_{1\leq i\leq n}|A_i||B_i|\geq n^{2-o(1)},\qquad
 |A_i+B_i|=|A_i||B_i|\quad(1\leq i\leq n),
\]
and
\[
 (A_i+B_i)\cap(A_j+B_k)=\varnothing
 \quad\text{for all }1\leq i,j,k\leq n\text{ with }j\ne k?
\]
More explicitly, for every $\epsilon>0$ and $M\in\mathbb N$, must some $n\geq M$ admit such subsets with $n^{2-\epsilon}\leq|H_n|\leq n^{2+\epsilon}$ and $|A_i||B_i|\geq n^{2-\epsilon}$ for every $i$, while the two exact addition conditions hold? \sref{167}{2}
\subsection{Short English statement}
Can a nearly quadratic-size abelian group support arbitrarily many nearly quadratic-size injective addition pairs whose matched sums avoid every mismatched sum?
\subsection{Sources}
\begin{description}
\item[\hypertarget{source:167:1}{S1}] ULAM.AI and the UnsolvedMath Contributors, \emph{UnsolvedMath: A Curated Collection of Open Mathematics Problems}, record 167 (GREEN-079). CC BY 4.0. \url{https://huggingface.co/datasets/ulamai/UnsolvedMath}.
\item[\hypertarget{source:167:2}{S2}] Henry Cohn, Robert Kleinberg, Balázs Szegedy and Christopher Umans, \emph{Group-theoretic algorithms for matrix multiplication}, FOCS 2005, 379--388; arXiv:math/0511460v1, §4, Definition 4.1 and Conjecture 4.7. Replacing each source subset $B_i$ by its negative converts the difference-set convention into the displayed sum-set convention. \url{https://arxiv.org/abs/math/0511460}.
\item[\hypertarget{source:167:3}{S3}] Ben Green, \emph{100 Open Problems}, maintained author list, Problem 36 and its comments, accessed 7 October 2026. \url{https://people.maths.ox.ac.uk/greenbj/papers/open-problems.pdf}.
\end{description}
\label{end:167}
\end{document}
